% chapters/04_results_part2.tex - 第4章 研究结果与分析（续）
% 4.9 对应目的4（用户接受度）
% 4.10 对应目的4（决策机制）
% 4.11 本章小结

\section{系统技术接受度量表结果}
\label{sec:acceptance}

AI辅助组11名受试者在T1后测阶段填写了系统可用性量表（SUS）与技术接受度量表（感知有用性PU、信任度Trust、持续使用意愿Intention）。个人评分与描述统计汇总见表~\ref{tab:acceptance}。

\begin{table}[htbp]
\centering
\caption{AI辅助组系统评价量表描述性统计与个人评分（n=11）}
\label{tab:acceptance}

\small
\begin{tabular}{lrrrr}
\toprule
编号 & SUS总分 & PU均分 & Trust均分 & Intention均分 \\
\midrule
P001 & 67.5 & 4.50 & 4.50 & 4.67 \\
P004 & 72.5 & 4.75 & 4.50 & 4.67 \\
P005 & 57.5 & 4.25 & 4.25 & 4.33 \\
P006 & 67.5 & 4.50 & 4.75 & 4.67 \\
P007 & 65.0 & 4.50 & 4.50 & 4.33 \\
P009 & 62.5 & 4.25 & 4.25 & 4.33 \\
P010 & 67.5 & 4.25 & 4.25 & 4.33 \\
P011 & 62.5 & 4.50 & 4.50 & 4.67 \\
P016 & 67.5 & 4.50 & 4.50 & 4.33 \\
P018 & 57.5 & 4.00 & 4.00 & 4.00 \\
P025 & 67.5 & 4.50 & 4.50 & 5.00 \\
\midrule
\textbf{均值$\pm$SD} & \textbf{65.00$\pm$4.61} & \textbf{4.41$\pm$0.20} & \textbf{4.41$\pm$0.20} & \textbf{4.49$\pm$0.27} \\
范围 & 57.5--72.5 & 4.00--4.75 & 4.00--4.75 & 4.00--5.00 \\
\bottomrule
\end{tabular}

\end{table}

SUS得分均值为65.00$\pm$4.61分（满分100分）。11名受试者中仅P004一人（9.1\%）SUS评分达到70分, 个人得分范围57.5--72.5分。感知有用性均分4.41$\pm$0.20, 信任度均分4.41$\pm$0.20, 持续使用意愿均分4.49$\pm$0.27, 三项TAM维度均分均超过4.4分（满分5分, 得分率88\%以上）。持续使用意愿最高者为P025（5.00分）, SUS、PU、Trust、Intention四项得分最低者均为P018（分别为57.5、4.00、4.00、4.00分）。

\section{质性访谈结果}
\label{sec:qualitative}

作为解释性序贯混合方法研究的质性阶段, 本研究在T1后测完成后对受试者开展了半结构化深度访谈, 从微观行为体验、心理感知与决策逻辑层面呈现两种模式下训练者的主观反馈内容。

\subsection{访谈实施与样本基本特征}
\label{sec:interview_sample}

访谈对象采用最大差异抽样与主题饱和相结合的策略确定。在完成T1量化测试后, 依据组别归属（AI辅助组vs.自我指导组）、基线力量分层（高力量层vs.低力量层）以及训练执行特征（是否触发自动减组、负荷自主调整频率）三个维度进行分层抽样, 共纳入12名受试者进行一对一深度访谈。其中, AI辅助组5名（P004、P010、P016、P018、P025）, 自我指导组7名（P003、P012、P013、P014、P015、P021、P027）。访谈对象的样本构成与特征分布见表~\ref{tab:interview_sample}。

实际纳入访谈的AI辅助组5名受试者均来自低力量层, 高力量层未覆盖；自我指导组则覆盖高、低力量层。

\begin{table}[htbp]
\centering
\caption{质性访谈受试者基本信息与抽样特征分布表（n=12）}
\label{tab:interview_sample}

\small
\begin{tabular}{llrcp{2cm}ccc}
\toprule
编号 & 组别 & 分层 & $\Delta$1RM & $\Delta$CMJ & 训练特征 & 时长 & 行为类型 \\
 \\
\midrule
P004 & AI & 低 & +15.0 & +6.7 & 命中率85\% & 25.4 & 积极型 \\
P010 & AI & 低 & +12.5 & +3.1 & 减组1次 & 27.2 & 安全型 \\
P016 & AI & 低 & +5.0 & $-$0.3 & 减组2次 & 28.0 & 安全型 \\
P018 & AI & 低 & +7.5 & +10.3 & SUS偏低 & 26.5 & 批判型 \\
P025 & AI & 低 & +22.5 & +6.7 & 命中率89\% & 24.8 & 积极型 \\
P003 & 自 & 低 & +7.5 & +3.8 & 调整2次 & 25.1 & 规则型 \\
P012 & 自 & 高 & +10.0 & +0.7 & 经验丰富 & 26.8 & 批判型 \\
P013 & 自 & 低 & +10.0 & +8.1 & 调整4次 & 27.5 & 规则型 \\
P014 & 自 & 高 & +10.0 & +1.5 & 保守力竭 & 25.9 & 批判型 \\
P015 & 自 & 高 & +17.5 & +0.6 & 慢速硬顶 & 26.2 & 负荷型 \\
P021 & 自 & 低 & +10.0 & +0.1 & RIR波动 & 24.5 & 负荷型 \\
P027 & 自 & 低 & +7.5 & +3.6 & 调整3次 & 26.0 & 规则型 \\
\bottomrule
\end{tabular}

\end{table}

访谈在T1测试结束后1至2天内于安静实验室环境进行, 单次访谈实际持续时间在24.5至28.0分钟之间, 平均26.2分钟。录音资料在征得受试者书面知情同意后完整采集, 转录形成12份中文逐字文本。分析遵循Braun与Clarke（2006）的反思性主题分析法, 并执行附录B设定的证据分层规则, 通过三轮编码提炼出6个核心主题与5类跨个案行为特征模式。

\subsection{核心主题呈现}
\label{sec:themes}

\textbf{主题一：结构化反馈作为负荷决策的外部参照}

两组受试者的叙事一致表明, 反馈机制的核心心理价值在于为自我训练场景中的负荷选择提供了外部参照。AI辅助组受试者普遍指出, 实时速度反馈为每一次动作执行提供了即时、客观的外部参照。P025与P004描述道, 在没有外部监控的自我训练中, 组间加重往往伴随犹豫；获得即时速度数据后, 只要上一组速度维持在目标区间中上限, 便能果断加重。在自我指导组中, 规范化RIR训练卡同样发挥了参照作用。P003与P027表示, 书面训练卡中明确的RIR目标改变了此前"要么力竭、要么随意"的状态, 引导其在动作末期主动评估自身的剩余能力。

\textbf{主题二：RIR边界的动作技术锚定与经验分化}

自我指导组7名受试者的深度访谈呈现了RIR概念在实际执行中的两种不同锚定方式。具备较长训练年限的受试者（P013、P014、P027）将RIR严格锚定在标准动作技术形态上。P013明确将"出现躯干过度前倾、膝内扣或动作轨迹明显停顿之前的重复"界定为有效储备次数；P027则根据当组后半程的动作速度衰减感知主动调控下组负荷或延长组间休息。训练经验有限的受试者（P015、P021）则呈现出不同的RIR认知方式。P015反映其在训练中往往将"咬牙慢速硬顶、伴随躯干代偿完成的最后几次"依然计入完成次数；P021反映在非力竭状态下难以准确体会剩余能力。

\textbf{主题三：自动减组触发后的心理体验演变}

在经历过自动减组的AI辅助组受试者中, 基于训练前纵跳准备度的动态减组规则触发引发了具有时间演变特征的心理体验。在触发减组的即时阶段, 受试者体验到心理落差。P010描述在系统提示"当日准备度不足、深蹲减少1组"的当下, 产生了"计划被打乱"与"自身状态未达标的挫败感"；P016在首次触发时曾产生对"训练刺激量被削减是否会影响力量增长"的担忧。在干预推进与后测回顾阶段, 受试者对自动减组的理解发生了变化。P010与P016均指出, 减组后当次训练的课后肌肉酸痛感明显轻于高负荷课次, 且在随后的下一次训练中展现出良好的力量状态。

\textbf{主题四：功能价值、操作摩擦与用户自主权诉求}

AI辅助组对数智化监控系统的接受度评价呈现出多层次结构。在功能价值层面, 受试者高度评价系统的功能实用性。在操作流程层面, 受试者反映了多维度的使用摩擦：P018强调每组试举前需要确认设备状态；P025指出当前专业设备形态更适合在专业实验室或双人搭档模式下使用。在用户自主权层面, P018明确表示, 虽然认可客观速度数据的参考价值, 但不希望系统成为绝对的"训练指挥官"；而P010、P025则表现出较强的指令遵从意愿。

\textbf{主题五：训练经验与工具功能诉求的对应关系}

跨个案叙事呈现出训练经验与数字化工具功能诉求之间的对应关系。低力量层中更倾向接受外部指令的受试者（P003、P010、P025）对数字化工具呈现出"行动指引与安全边界"诉求, 希望系统提供直接、明确的操作指令。高力量层受试者（P012、P014）以及低力量层但训练经验丰富、自主性诉求较强的受试者（P018）则呈现出"数据验证与疲劳诊断"诉求, 希望借助高精度的速度与位移数据验证自身的本体感觉, 并保留最终的负荷决定权。

\textbf{主题六：体能适应、系统评价与安全体验的独立性}

访谈资料与量化结局的对照呈现出：训练者的体能提升幅度、对系统的可用性评价以及主观安全感之间不存在简单的对应关系。P018在4周干预中取得了全样本最高的CMJ增量（+10.3 cm）, 同时在访谈中对系统的操作流程与智能化体验提出了尖锐批判；P025取得了全样本最高的深蹲1RM增量（+22.5 kg）, 同时展现出高系统信任度与使用意愿；P016的体能增量相对平缓（1RM +5.0 kg, CMJ $-$0.3 cm）, 但对算法安全保护给予积极评价；自我指导组中严格维持动作质量的P013取得了双重适应（1RM +10.0 kg, CMJ +8.1 cm）。

\subsection{跨个案行为类型归纳}
\label{sec:behavior_patterns}

基于12名访谈对象的决策风格、规则遵从度、自主权诉求与体验特征, 归纳出5种典型的抗阻训练行为特征模式, 见表~\ref{tab:behavior_patterns}。

\begin{table}[htbp]
\centering
\caption{抗阻训练行为特征跨个案类型归纳表（n=12）}
\label{tab:behavior_patterns}

\small
\begin{tabular}{lp{2cm}p{8cm}}
\toprule
类型命名 & 代表个案 & 核心行为特征 \\
\midrule
积极接受型 & P004, P025 & 严格遵从速度反馈与建议；将数据作为加重与突破信心的客观支撑 \\
安全依赖型 & P010, P016 & 极度关注疲劳与损伤风险；高度认可准备度减组与安全停组机制 \\
批判接受型 & P018, P012, P014 & 具备扎实的自主训练经验；高度重视本体感觉；批判操作繁琐与算法僵化 \\
规则内化型 & P003, P013, P027 & 深刻理解RIR逻辑；将动作技术变形与代偿作为RIR归零的刚性边界 \\
负荷导向型 & P015, P021 & 偏好追求大重量与极限重复；容易将代偿硬顶计入有效储备 \\
\bottomrule
\end{tabular}

\end{table}

\section{本章小结}
\label{sec:ch4_summary}

本章按研究目的1至4的展开顺序系统呈现了量化实验与质性访谈的全部研究结果。核心数据事实归纳如下。

\textbf{回应目的1（工具效度基础）}

第一, 研究一在受控实验条件下证实SportSci Pro与GymAware在杠铃后蹲MCV测量上达到优秀一致性（ICC=0.940, MAE=0.047 m/s, Bias=+0.002 m/s）。

第二, 研究二真实训练场景配对分析显示, Rep级（n=43, 3名受试者）ICC=0.991（Bias=+0.0019 m/s）, 热身级（n=88, AI辅助组全体11名受试者）ICC=0.987, 但存在+0.0212 m/s的稳定系统性正向偏移（88/88对均为正偏）。

\textbf{回应目的2（干预可行性）}

第三, 研究二24名PP受试者的课次层训练完成率达98.4\%（189/192）, 高依从率达100.0\%（24/24）, AI辅助组速度目标命中率达80.2\%$\pm$6.9\%, 自动减组触发忠实度达100\%, 严重不良事件为0例。9项预设推进标准中7项达标, 主要后测完成率（66.7\%, 24/36）与AI辅助组SUS评分（65.00$\pm$4.61分）2项未达标。

\textbf{回应目的3（干预效果）}

第四, 两组在全期总外部负荷（17651.8 vs. 17461.5 kg, $p$=0.854）、深蹲主项负荷（$p$=0.988）、全期平均sRPE（6.5 vs. 6.4）及个体Hooper均值（20.4 vs. 19.8）上高度对等；在总完成重复次数上, AI辅助组显著少于自我指导组（100.3 vs. 105.5次, $p$=0.035）, 派生的每次重复平均负荷AI辅助组高出约6.1\%。

第五, Hooper LMM分析（Kenward-Roger）显示组别主效应$p$=0.421、交互效应$p$=0.945、课次主效应$p$=0.004；两组Hooper均值在时间演变上均呈先升后降形态（AI辅助组S1=19.91$\to$S4峰值21.45$\to$S8=18.45；自我指导组S1=18.54$\to$S6峰值20.85$\to$S8=18.62）。

第六, ANCOVA分析显示两组在深蹲绝对1RM（调整后差值+2.43 kg, $p$=0.266, $g$=0.46）、相对1RM（+0.03 kg/kg, $p$=0.273, $g$=0.39）与SJ高度（+0.39 cm, $p$=0.622, $g$=0.43）上未达统计学显著；在CMJ高度上调整后差值+2.16 cm（$p$=0.051, $g$=0.85, $g$ 95\% CI：0.02至1.69）, FDR校正后$p$=0.126；在训练自我效能上调整后差值+9.88分（$p<0.001$, $g$=2.70, FDR校正后仍显著）。

第七, 训练负荷与结局变化的28对Spearman探索性相关中, 8对达到原始$p<0.05$, FDR校正后全部不显著。

\textbf{回应目的4（用户接受度与决策机制）}

第八, AI辅助组SUS得分65.00$\pm$4.61分, 11人中1人达70分；感知有用性、信任度、使用意愿三项TAM维度均分分别为4.41、4.41、4.49（满分5分）。

第九, 12名受试者的半结构化访谈通过反思性主题分析形成6个核心主题与5类跨个案行为特征模式（积极接受型、安全依赖型、批判接受型、规则内化型、负荷导向型）。
